% Редакційне доопрацювання; числові свідчення залишаються в заморожених результатах.
\section{Вступ}\label{sec:intro}

\subsection{Мотивація}
Згаданий дослідницький матеріал може бути нотаткою про антену, абзацом у звіті, таблицею, файлом вихідного коду чи реченням у записі. Його формулювання може відрізнятися від запиту, а мова документа --- від мови пошуку. Архів також може містити неопубліковані матеріали й розмови з агентами, які власник воліє обробляти локально. Тому лексичний збіг, семантична близькість, походження джерела й обчислювальні витрати стають частинами одного пошукового завдання.

Ми досліджуємо архіви, у яких українська й російська мови трапляються разом. Флексія створює практичний вибір у лексичному пошуку: зіставлення словоформ зберігає написаний вираз, тоді як розширення лемами може пов'язувати різні форми слова. Векторні вкладення дають інший зв'язок між запитом і документом, зокрема між мовами та модальностями. Поєднання цих сигналів потребує ще одного рішення: зливати ранжовані списки чи обирати режим пошуку за лексичними ознаками. Ефективність кожного рішення для конкретного архіву потребує перевірки.

archive-index реалізує ці альтернативи в спільному локальному конвеєрі індексації та пошуку. Система поєднує SQLite FTS5/BM25, морфологічний аналіз української й російської, представлення EmbeddingGemma~2 та налаштовуване гібридне ранжування (розділ~\ref{sec:system}). Наше питання: \emph{які компроміси між якістю пошуку й ресурсними витратами демонструють ці реалізовані конфігурації на побудованому багатомовному архіві та наскільки відповідь залежить від модельних оцінок релевантності?}

\subsection{Внесок}
Дослідження надає чотири групи доступних для перевірки матеріалів і порівнянь:
\begin{enumerate}[label=(\alph*)]
\item \emph{Реалізована система.} Інкрементальна індексація дванадцяти колекцій джерел, локатори фрагментів, лексичне розширення та спільні представлення тексту, зображень і мовлення з явними перемикачами порівнюваних конфігурацій.
\item \emph{Побудований бенчмарк.} \uapar{} містить 7\,070 файлів із публічних джерел, 200 модельних запитів і 5\,572 пари «запит--документ» у пулі, оцінені чотирма модельними суддями. Окреме розширення з 50 записів додає цілі серед документів, таблиць і файлів репозиторіїв; воно містить 49 різних формулювань. Визначення запитів, прогони та всі початкові оцінки збережено.
\item \emph{Порівняння з аналізом чутливості до оцінювання.} Десять конфігурацій та абляцій порівнюються з оцінками невизначеності за запитами, парними тестами й кількома варіантами міток. Контролі на склад кандидатів і виправлення покриття відділяють відмінності енкодерів від змін пошукового конвеєра.
\item \emph{Опис витрат і роботи в конкретному середовищі.} Вимірювання CPU/GPU на одному ноутбуці, порівняння пошуку зі збірками GGUF, синтетичне масштабування точного пошуку/Faiss/HNSW та агреговані перевірки приватного архіву автора виявляють різні межі реалізації й оцінювання.
\end{enumerate}
Система поєднує відомі пошукові компоненти. Емпіричний внесок --- їхнє відтворюване порівняння в цих архівних умовах, зокрема результати, що не підтримують початкової гіпотези про адаптивне злиття. Додаток~\ref{sec:intro-history} зберігає виправлення реалізації й зміни протоколу, відрізняючи зміни до заморожених прогонів від пізнішого контролю покриття.

\subsection{Свідчення та межі висновків}
Ми розрізняємо чотири рівні свідчень:
(L1)~модульні й інтеграційні перевірки поведінки реалізації;
(L2)~текстовий пошук на \uapar{} з модельними оцінками (розділи~\ref{sec:uapar}--\ref{sec:results});
(L3)~операційні перевірки архіву автора, експортовані лише агрегатами (розділ~\ref{sec:field});
(L4)~вимірювання затримки, пам'яті та використання CPU/GPU на одному ноутбуці.
Окремі завдання із зображеннями й мовленням використовують пари «підпис--зображення» та «речення--запис» із Crossmodal-3600 і FLEURS. Вони не дають людських оцінок релевантності текстового бенчмарку.

Ця версія не містить людського аудиту релевантності. Оцінки текстового пошуку залежать від модельних міток і побудованих запитів; усі 170 початкових цілей запитів із відповіддю --- нотатки. Окреме розширення додає формати цілей, проте не робить жодного набору репрезентативним для особистих інформаційних потреб. Дослідження не навчає моделей, не оцінює інших мов, крім української, російської й наявної в коді англійської, а для масштабування понад публічну колекцію використовує синтетичні вектори. Тести реалізації та перевірки приватності не сертифікують усіх аспектів коректності чи конфіденційності.
